\section{An arbitrary-input-law oracle for bounded monotone activations}
\label{sec:all-d-product-midpoint}
The random-subset midpoint construction also gives a simpler oracle in
arbitrary input dimension. It uses the bounded ridge-feature class rather
than a spatial smoothness assumption on the input distribution.

Fix a bounded, nonconstant, nondecreasing, globally Lipschitz activation
$\phi:\mathbb R\to\mathbb R$. Tanh, logistic sigmoid, and hardtanh are
examples. Constants below may depend on its fixed supremum $K$ and
Lipschitz constant. Let $X$ have any probability law on $[-1,1]^d$ and
let $|Y|\le1$ and $d\ge1$. Put $D=d+1$ and $\widetilde x=(x,1)$.
Let $a_0=d$ for the original first-layer convention; the alternative
$a_0=d+1$ is also covered. The learner is
\[
 f_w(x)=N^{-1}\sum_{j=1}^N G_j\phi(\theta_j\cdot\widetilde x),
 \qquad
 \theta_j\sim N(0,\nu I_D/a_0),\quad G_j\sim N(0,\nu),
 \quad0<\nu\le1,
\]
with all coordinates independent. Thus the actual readout coefficients
$G_j/N$ have variance $\nu/N^2$. Define
\[
 Q_G=N^{-1}\sum_jG_j^2,\qquad \omega=N/\nu,\qquad
 d\Pi_\lambda=Z_\lambda^{-1}
 e^{-\omega\widehat R(f_w)/(2\lambda)}d\mu.
\]
Here $R(f)=\mathbb E(f(X)-Y)^2$ and
$\widehat R(f)=n^{-1}\sum_i(f(X_i)-Y_i)^2$.
The physical Gaussian observation variance is $n\lambda\nu/N$.
Both the likelihood and the returned mean use the unrestricted raw output.

Let a deterministic comparison network be
\[
 g(x)=\sum_{i=1}^C c_i\phi(\vartheta_i\cdot\widetilde x),
 \qquad \|g\|_\infty\le1.
\]
The supremum is over the input cube. If $V=\sum_i|c_i|>0$, set
$p_i=|c_i|/V$, omitting zero coefficients, and define
\[
 J_g=a_0\sum_i p_i\|\vartheta_i\|_2^2,
 \qquad
 A_n(g)=1+V^2+J_g+D\log(e n(V+1)).                       \tag{AD1}
\]
For $g=0$ use $V=J_g=0$. The term $J_g$ is measured in the actual
fan-in-normalized incoming Gaussian energy. In particular, if every
teacher coefficient has absolute value at most $B\ge1$, then
\[
 A_n(g)\le C_\phi\{1+C^2B^2+D^2B^2+D\log(e n(CB+1))\}. \tag{AD2}
\]
An incoming row-norm cap $B_{\rm in}$ instead gives
$J_g\le a_0 B_{\rm in}^2$, independently of the teacher width.

\begin{lemma}[Capacity of the bounded ridge hull]
\label{lem:all-d-ridge-hull}
Let $\mathcal F_q$ be the class of all finite sums
$\sum_j v_j\phi(\theta_j\cdot\widetilde x)$ with
$\sum_j|v_j|\le\sqrt q$. On every list of $n$ inputs,
\[
 \operatorname{Rad}_n(\mathcal F_q)
 \le C_\phi\sqrt{Dq/n},\qquad
 \sup_{f\in\mathcal F_q}\|f\|_\infty\le K\sqrt q.        \tag{AD3}
\]
Every learner output belongs to $\mathcal F_{Q_G}$, without any bound
on incoming weights.
\end{lemma}
\begin{proof}
For a monotone activation, each subgraph test
$t<\phi(\theta\cdot\widetilde x)$ is a constant test or an affine
halfspace test in $(x,\phi^{-1}(t))$, with a generalized inverse at
flat portions. Strict or weak boundary choices give the same VC bound.
The ridge class therefore has VC-subgraph dimension at most $D+1$.
The standard VC-subgraph covering bound for a bounded real-valued class
is $N(\epsilon K,\mathcal H,L^2(P_n))\le(A/\epsilon)^{cD}$,
with universal $A,c$. Dudley's entropy integral gives worst-case
empirical Rademacher complexity $C K\sqrt{D/n}$. Adding negatives
only changes the constant. A signed convex hull does not increase
Rademacher complexity beyond its radius, proving (AD3).
Finally $N^{-1}\sum_j|G_j|\le\sqrt{Q_G}$.
\end{proof}

\begin{lemma}[A common relative-loss event]
\label{lem:all-d-relative}
Fix a deterministic $u$ with $\|u\|_\infty\le2$. There is a constant
$C_\phi$ such that, for $n\ge2$, putting
\[
 b_n=\frac{C_\phi(D+1)\log^4(e n/\delta)}n,              \tag{AD4}
\]
with probability at least $1-\delta$, simultaneously for all
$q\le n^{22}$ and $f\in\mathcal F_q$,
\begin{align}
 R(f)-R(u)\le{}&\widehat R(f)-\widehat R(u)
 +\tfrac16\|f-u\|_n^2+c\lambda(q+1)+C_{\phi,c}b_n/\lambda,
                                                                  \tag{AD5}
\end{align}
whenever $n^{-2}\le\lambda\le1$ and $\lambda\ge C_{\phi,c}b_n$.
The positive constant $c$ may be chosen arbitrarily small by enlarging
the constants in the likelihood floor and final remainder.
\end{lemma}
\begin{proof}
Write $h=f-u$. Apply the standard smooth-loss Rademacher inequality
for nonnegative squared loss (Srebro, Sridharan, and Tewari, 2010,
Theorem~1~\cite{sst}) to the class
$\mathcal F_q-u$, with response zero. Its Rademacher complexity is
at most $C_\phi\sqrt{D(q+1)/n}$ and its squared envelope is at most
$C_\phi(q+1)$. For any fixed small $\eta>0$, Young's inequality gives
\[
 (P_X-P_n)h^2\le\eta P_nh^2+C_{\phi,\eta}b_n(q+1).      \tag{AD6}
\]
Only this one-sided norm comparison is needed. The invoked smooth-loss
bound is the usual inequality
\begin{align*}
 P h^2-P_nh^2\le{}&C\sqrt{P_nh^2}
 \left[\log^{3/2}(2n)\operatorname{Rad}_n(\mathcal H)
       +\sqrt{B\log(1/\zeta)/n}\right]\\
 &+C\log^3(2n)\operatorname{Rad}_n(\mathcal H)^2
 +CB\log(1/\zeta)/n.
\end{align*}
where $|h|^2\le B$. Its hypotheses are satisfied by the bounded class
in (AD3).

For the other term, the multiplier $u(X)-Y$ is bounded by three.
Symmetrization, contraction, (AD3), and bounded-difference concentration
give
\[
 |(P-P_n)((u-Y)h)|\le C_\phi\sqrt{b_n(q+1)}.             \tag{AD7}
\]
This argument does not require the multiplier to be independent of $X$:
contraction is applied conditionally to the observed pairs $(X_i,Y_i)$.
Apply both inequalities at dyadic radii through $2n^{22}$, including
radius one, and allocate confidence equally to these $O(\log n)$
events. This explains the safe logarithmic factor in (AD4), and gives
one event for all intervening radii.

Use the exact identity
\[
 R(f)-R(u)-\widehat R(f)+\widehat R(u)
 =(P_X-P_n)h^2+2(P-P_n)((u-Y)h).
\]
Take $\eta=1/6$. The term $C b_nq$ is absorbed into a prescribed
small multiple of $\lambda(q+1)$ when $\lambda\ge Cb_n$.
Young's inequality bounds $C\sqrt{b_n(q+1)}$ by another such multiple
plus $C b_n/\lambda$.
\end{proof}

\begin{lemma}[Variation-weighted comparison packs]
\label{lem:all-d-comparison}
There is a one-pack comparison law $q_g$ such that its mean function
$g_\sigma$ is independent of $\nu$,
\[
 \|g_\sigma-g\|_\infty\le C_\phi n^{-20},\qquad
 \mathbb E_{q_g}G^2\le V^2+1,
 \qquad
 \nu\operatorname{KL}(q_g\Vert\mu_0)\le C_\phi A_n(g).  \tag{AD8}
\]
Consequently, whenever $N\lambda\ge C_\phi$,
\[
 \Pi_\lambda[\widehat R(f)-\widehat R(g_\sigma)]
 +\tfrac13\Pi_\lambda\|f-g_\sigma\|_n^2
 +c_\phi\lambda\Pi_\lambda Q_G
 \le C_\phi\lambda A_n(g).                             \tag{AD9}
\]
\end{lemma}
\begin{proof}
For $V>0$, choose index $i$ with probability $p_i$, then independently
\[
 \theta\sim N(\vartheta_i,\sigma^2I_D/a_0),\qquad
 G\sim N(V\operatorname{sign}c_i,\nu),\qquad
 \sigma=\frac{n^{-20}}{V+1}.
\]
Since $\|\widetilde x\|_2\le\sqrt D$, the incoming perturbation
changes a ridge in mean by at most
$\sqrt{D/a_0}\operatorname{Lip}(\phi)\sigma
 \le\sqrt2\operatorname{Lip}(\phi)\sigma$.
Thus the mean error is at most $C_\phi V\sigma$.
The readout second moment is exactly $V^2+\nu$.
Convexity and the Gaussian entropy formula give
\[
 \nu\operatorname{KL}(q_g\Vert\mu_0)
 \le\tfrac12\{J_g+V^2+D\sigma^2-D\nu
                     -D\nu\log(\sigma^2/\nu)\},
\]
which is bounded by $C_\phi A_n$ because $\nu\log\nu\le0$.
For $g=0$ simply take the native centered pack law.

Apply the random-subset replacement construction of
Lemma~\ref{lem:tanh-subset-midpoint}: keeping a uniformly
chosen $k=\lfloor N/2\rfloor$ posterior packs and replacing the others
by independent $q_g$ draws contracts relative entropy by $p=k/N$.
Its conditional variance is bounded by
\[
 K^2\left\{\frac{p(1-p)}{N-1}Q_G
             +\frac{1-p}{N}\mathbb E_{q_g}G^2\right\}.
\]
Gibbs optimality, the squared-loss midpoint identity, and division by
$1-p$ give
\begin{align*}
 &\Pi_\lambda[\widehat R(f)-\widehat R(g_\sigma)]
 +p\Pi_\lambda\|f-g_\sigma\|_n^2
 +\frac{2\lambda\nu}{N}\operatorname{KL}(\Pi_\lambda\Vert\mu)\\
 &\hspace{8mm}\le2\lambda\nu\operatorname{KL}(q_g\Vert\mu_0)
 +K^2\left\{\frac p{N-1}\Pi_\lambda Q_G
                  +\frac{\mathbb E_{q_g}G^2}N\right\}.
\end{align*}
Use only the readout-block Gaussian moment:
\[
 \operatorname{KL}(\Pi_\lambda\Vert\mu)
 \ge\frac\omega4\Pi_\lambda Q_G-\frac N2\log2.
\]
Absorb the $Q_G/N$ term using $N\lambda\ge C_\phi$, and use
$(V^2+1)/N\le C_\phi\lambda A_n$. This proves (AD9).
\end{proof}

\begin{theorem}[A sharp arbitrary-input-law Gaussian oracle]
\label{thm:all-d-sharp}
For every fixed comparison network $g$ above, with probability at least
$1-\delta$, simultaneously for all $0<\nu\le1$, all
$n^{-2}\le\lambda\le1$ satisfying $\lambda\ge C_\phi b_n$, and all
$N\lambda\ge C_\phi$,
\[
 R(\mathbb E_{\Pi_\lambda}f)\le\Pi_\lambda R(f)
 \le R(g)+C_\phi\left\{\lambda A_n(g)+\frac{b_n}{\lambda}\right\}
       +C_\phi n^{-10}.                                \tag{AD10}
\]
There is no density or anti-concentration condition on the input law.
No output truncation, prior restriction, or architecture mixture is used.
\end{theorem}
\begin{proof}
The prior empirical risk is at most $M=1+K^2$, since the readout is
centered and $\mathbb E_\mu f(x)^2\le K^2\nu/N$.
Thus Jensen gives $Z_\lambda\ge\exp[-\omega M/(2\lambda)]$.
The law of $\omega Q_G$ is $\chi^2_N$, so above a fixed constant
$\mu(Q_G>t)\le\exp(-c\omega t)$. Integrating the corresponding
posterior bound shows that the omitted risk, empirical loss and
quadratic-distance integrals on $Q_G>n^{22}$ are at most $C_\phi n^{-10}$,
uniformly for $\lambda\ge n^{-2}$, $N\ge1$ and $0<\nu\le1$.
For the finitely many smaller $n$, increase the deterministic lower
sample-size threshold by an activation-dependent constant.
All these quantities are bounded by $C_\phi(1+Q_G)$.

Use (AD5) with $u=g_\sigma$ and a constant $c$ smaller than half
$c_\phi$ in (AD9). Its norm is at most two for sufficiently large $n$.
Integrate over the posterior, restore the negligible tail, and apply
(AD9). The extra constant $c\lambda$ is absorbed by $C\lambda A_n$
since $A_n\ge1$. The remaining distance and readout-energy terms are nonnegative
and can be dropped. Finally $R(g_\sigma)\le R(g)+C_\phi n^{-20}$.
Jensen proves the mean-risk statement. The sample event is fixed by
$g_\sigma$ and the width/scale-independent ridge classes; $g_\sigma$
does not depend on $\nu$, so the advertised simultaneity is valid.
\end{proof}

\begin{corollary}[An explicit polynomial]
\label{cor:all-d-sharp-polynomial}
In the nontrivial range $A_n(g)b_n\le c_\phi$, choose
\[
 \lambda=C_\phi\sqrt{b_n/A_n(g)},\qquad
 N\ge C_\phi\sqrt{A_n(g)/b_n}.
\]
Then
\[
 R(\mathbb E_{\Pi_\lambda}f)\le\Pi_\lambda R(f)
 \le R(g)+C_\phi\sqrt{A_n(g)b_n}+C_\phi n^{-10}.         \tag{AD11}
\]
The scale uses a declared deterministic upper bound for $A_n$, not the
unknown risk or a prior over teacher sizes. At fixed $d,B$ this is
$\widetilde O(C/\sqrt n)$ against width-$C$ teachers, and a sufficient
sample size for excess $\varepsilon$ is
\[
 n\ge\widetilde O\left(
 (d+1)[C^2B^2+(d+1)^2B^2]\varepsilon^{-2}\right).        \tag{AD12}
\]
For bounded readout variation and bounded incoming row norm, the resource
in (AD1) is independent of the number of teacher neurons. The comparator is still required to have
bounded raw output. The same risk bound holds against the infimum over
a declared such class by selecting a deterministic population-risk
near-minimizer; no common event over all data-selected comparisons is
asserted. In the nontrivial range, every $N\ge C_\phi n$ satisfies the
width requirement, after enlarging constants and logarithmic factors.
\end{corollary}
\begin{proof}
Balance the two terms of (AD10). The nontrivial-range condition makes
$\lambda\ge C_\phi b_n$, after the constants are chosen in order, and
also $\lambda\le1$ when $b_n$ is sufficiently small. Equation (AD2)
gives (AD12). Finally $A_nb_n\le c$ implies
$\sqrt{A_n/b_n}\le C/b_n\le Cn$.
\end{proof}

\paragraph{An elementary bound without an output cap.}
Theorem~\ref{thm:uncapped-all-d} below removes the output cap without
worsening (AD11). The following simpler envelope-based argument gives
a coarser statement for any finite comparison network.
Let $F=\|g\|_\infty$, so $F\le K V$. In the nontrivial range, the same
argument gives
\[
 R(\mathbb E_{\Pi_\lambda}f)\le\Pi_\lambda R(f)
 \le R(g)+C_\phi\left\{\lambda A_n(g)
             +(F+1)^2b_n/\lambda\right\}+C_\phi(F+1)n^{-10}.
                                                               \tag{AD13}
\]
One may assume $F\le n^2$ to keep the same polynomial energy cutoff;
the nontrivial optimized range below automatically has this property.
Indeed the comparison mean is bounded by $F+1$. The squared-norm
argument now has envelope $C_\phi(q+(F+1)^2)$, and its remainder is
$C_\phi b_n(q+(F+1)^2)$. The bounded multiplier argument gives
$C_\phi(F+1)\sqrt{b_n(q+(F+1)^2)}$. Young's inequality bounds these by
a small multiple of $\lambda q$, plus
$C_\phi\lambda(F+1)^2+C_\phi(F+1)^2b_n/\lambda$.
Since $(F+1)^2\le C_\phi A_n$, the new constant penalty is absorbed
in $C_\phi\lambda A_n$; the midpoint comparison is unchanged.
The loss-smoothing error is $C_\phi(F+1)n^{-20}$ and the Gaussian
tail remains negligible at the stated polynomial cutoff.

Choosing $\lambda=C_\phi(F+1)\sqrt{b_n/A_n}$ gives excess
$\widetilde O_\phi((F+1)\sqrt{(d+1)A_n/n})$, provided this is
smaller than a sufficiently small fixed constant and
$N\lambda\ge C_\phi$. With entrywise coefficient cap $B$ and fixed
$d,B$, this is $\widetilde O(C^2/\sqrt n)$ without a comparison-output
restriction. It is weaker than (AD11), whose bounded-output premise
is what gives the better width polynomial.

An explicit teacher readout bias can be included as one extra constant
neuron: choose a fixed $t_0$ with $\phi(t_0)\ne0$, take incoming row
$(0,\ldots,0,t_0)$, and coefficient $b_{\rm out}/\phi(t_0)$.
It is counted in the variation and incoming-energy resources above.

\paragraph{Marginally sub-Gaussian labels.}
The bounded-label hypothesis can be replaced by the marginal condition
$\mathbb E\exp(Y^2/K_Y^2)\le2$ for a fixed $K_Y<\infty$.
This does not require a bounded conditional mean or a conditionally
sub-Gaussian residual. Keep the bounded-output comparator of (AD1), and set
\[
 B_n=\max\{1,K_Y\sqrt{40\log(e n/\delta)}\},\qquad
 b_n^{\rm sg}=C_\phi(B_n+1)^2(D+1)\log^4(e n/\delta)/n.
\]
Equations (AD10)--(AD11) hold with $b_n^{\rm sg}$ in place of $b_n$
and an $O_{\phi,K_Y}(n^{-10})$ negligible term. The comparison pack and
resource $A_n$ are unchanged. In particular the optimized excess is
\[
 \widetilde O_{\phi,K_Y}\bigl(\sqrt{(d+1)A_n/n}\bigr),               \tag{AD14}
\]
with the same condition $N\lambda\ge C_\phi$. As before, the scale must
satisfy $n^{-2}\le\lambda\le1$ and
$\lambda\ge C_\phi b_n^{\rm sg}$; the optimized statement concerns its
nontrivial range. No observed label is modified in the actual algorithm.

Here are the tail and risk details. In the proof alone define
$Y^{B_n}=\operatorname{clip}(Y,-B_n,B_n)$. With probability at least
$1-\delta/4$, every observed label equals its clipped value. The pairs
$(X,Y^{B_n})$ retain exactly the original input law. Apply the preceding
common statistical event to this clipped-label distribution. Its
multiplier estimate gains a factor $B_n+1$, which changes $b_n$ to
$b_n^{\rm sg}$. The midpoint identity and comparison-pack entropy do
not require bounded labels. On the no-exceedance event the actual
posterior, including its normalizer, is exactly the posterior appearing
in that analysis, and its native empirical risk is at most $B_n^2+K^2$.
The Gaussian readout-tail proof consequently gives
\[
 \Pi_\lambda Q_G\le C_\phi\{1+(B_n^2+1)/\lambda\}.
\]
The eligibility condition $b_n^{\rm sg}\lesssim1$ implies $B_n^2\lesssim n$;
thus $Q_G>n^{22}$ still has negligible posterior integrals uniformly
for $\lambda\ge n^{-2}$.

Most importantly, only a relative-risk transfer is needed:
\begin{align*}
 &[R_Y(f)-R_Y(g)]-[R_{Y^{B_n}}(f)-R_{Y^{B_n}}(g)]\\
 &\hspace{15mm}=-2\mathbb E[(Y-Y^{B_n})(f(X)-g(X))].
\end{align*}
The $Y^2$ terms cancel exactly. The marginal tail assumption gives
$\mathbb E|Y-Y^{B_n}|\le C K_Y(e n/\delta)^{-20}$.
Since $\|f\|_\infty\le K\sqrt{Q_G}$ and $\|g\|_\infty\le1$,
posterior averaging and the displayed energy bound make the absolute
transfer error $O_{\phi,K_Y}(n^{-10})$. A confidence allocation between
this no-exceedance event and the clipped-label statistical event proves
(AD14), with exact leading comparator coefficient one and a common
event over the eligible widths and scales.

\paragraph{A hardtanh learner competing with shallow ReLU teachers.}
The hardtanh instance also gives an exact, same-depth comparison with
coefficient-bounded ReLU teachers. Write
\begin{gather*}
 g(x)=b_{\rm out}+\sum_{i=1}^C c_i\operatorname{ReLU}(a_i\cdot x+b_i),\\
 |b_{\rm out}|,|c_i|,|(a_i)_j|,|b_i|\le B,
 \qquad B\ge1,\qquad M=(d+1)B.
\end{gather*}
On the input cube, $|a_i\cdot x+b_i|\le M$, and the exact identity
\[
 \operatorname{ReLU}(t)
 =\frac M2\{\operatorname{hardtanh}(2t/M-1)+1\},\qquad |t|\le M,
\]
converts $g$ to a hardtanh teacher of width at most $C+1$.
The extra neuron is the constant $\operatorname{hardtanh}(1)=1$.
Its readout variation obeys
\[
 V'\le M\sum_i|c_i|+|b_{\rm out}|
 \le(d+1)CB^2+B.
\]
Each transformed incoming row has Euclidean norm at most three,
including the constant row, so $J'\le9a_0$.
The comparison function itself is unchanged. Therefore (AD11) applies
when $\|g\|_\infty\le1$, giving fixed-$d,B$ sample complexity
$\widetilde O(C^2\varepsilon^{-2})$ under any input law; (AD13) gives
the coarser no-output-cap statement. Marginally sub-Gaussian labels are
covered by (AD14) when the comparison output is bounded.
The actual learner here has hardtanh activation. This is not a claim
that the raw ReLU learner has the bounded-ridge capacity bound (AD3).

